% ======================================================================
% SECTION 1: INTRODUCTION
% File: sections/introduction.tex
% ======================================================================

The integration of Physical AI systems into oncology clinical trials has established a new
paradigm for cancer treatment evaluation, delivery, and monitoring. Prior work in this
repository has developed the regulatory foundations (adapted 21 CFR Part 312 \cite{cfr312-adapt},
21 CFR Part 50 \cite{cfr50-adapt}, and ICH E6(R3) \cite{ich-adapt}), quantitative readiness
standards (Physical AI Standard Level and Unification Standard Level scoring
\cite{usl-standard}), national MCP server infrastructure \cite{national-mcp-paper}, federated
learning frameworks \cite{fl-paper}, site establishment documentation \cite{site-docs}, and a
complete single-patient journey orchestration \cite{patient-journey-paper}. These developments
address the trial execution layer but leave the sponsor layer, the organizational entity that
designs, funds, governs, and takes accountability for the trial, operating under a traditional
human-staffed model.

Traditional oncology trial sponsors employ hundreds of staff distributed across ten or more
functional areas: program leadership, clinical science and strategy, clinical operations,
medical monitoring and safety, data management and clinical data science, biostatistics,
regulatory affairs and operations, quality assurance and compliance, clinical supply chain,
and medical affairs and publications \cite{ich-e6r3}. This organizational model creates
bottlenecks at every decision gate. Phase I trials cost a median of \$14.2 million, Phase II
trials \$25.5 million, and Phase III trials \$65.8 million (Tufts CSDD, 2022)
\cite{tufts-csdd-cost}, with each day of Phase III delay costing \$55,716 in opportunity
cost \cite{tufts-delay}. The median development timeline for innovative drugs spans 8.3 years
\cite{tufts-timeline}, during which human coordination across sponsor functions introduces
delays at study startup, enrollment, data management, and regulatory submission.

Physical AI oncology trials demand a sponsor operating model that matches the autonomous
capabilities of robotic execution systems. When a surgical robot performs a tumor resection,
a collaborative robot prepares chemotherapy doses, and a motion-tracking robot positions
patients for radiation therapy, the sponsor oversight system must operate at corresponding
speed and precision. This paper presents a fully autonomous AI-native sponsor operating system
that replaces each human sponsor function with a software agent or robotic subsystem. The
system is not merely an electronic data capture platform, a clinical trial management system,
a robot controller, a protocol copilot, a safety dashboard, or a site orchestration tool. It
is a company-shaped trial execution application: an integrated operating system that performs
every function a pharmaceutical sponsor organization performs, from portfolio strategy through
regulatory submission and 25-year archiving.

The autonomous sponsor architecture comprises twelve functional agents organized into four
layers. The Governance Layer includes the portfolio agent (program prioritization and
investment logic), asset lead agent (asset strategy and vendor coordination), and board agent
cluster (go/no-go decision algorithms with escalation policies). The Study Execution Layer
encompasses the clinical accountability agent (protocol medical logic and dose optimization),
study orchestrator (cross-functional coordination and dependency management), clinical
operations agent (site startup, monitoring, and enrollment), safety agent (adverse event
intake, signal detection, and expedited reporting), regulatory agent (IND lifecycle management
and multi-jurisdiction filing), quality agent (SOP enforcement, CAPA lifecycle, and RBQM),
supply agent (IMP logistics, CMC lifecycle, and robotic pharmacy interface), and data and
biostatistics agent (data pipeline automation and CDISC compliance). The Site and Robotics
Layer includes the site gateway (site qualification, training, and investigator coordination)
and robot execution gateway (task-order generation, safety gate evaluation, and procedure
state management). The Trust Layer provides identity management, immutable audit
infrastructure, provenance chain tracking, and policy enforcement across all agent operations.

The system is not a replacement for the legal sponsor of record. Under 21 CFR \S312.50,
the sponsor bears responsibility for the overall conduct of the clinical investigation
\cite{cfr-part-312}. Under ICH E6(R3), the sponsor must implement a quality management system
proportionate to trial risks \cite{ich-e6r3}. The autonomous sponsor operates as an AI-native
operating system wrapped by a human or legal entity that retains accountability, override
authority for safety-critical decisions, and regulatory responsibility. When agent confidence
drops below defined thresholds or when mandatory human approval gates are triggered (such as
for high-risk robotic procedures or serious adverse event escalation), the system escalates to
the sponsor-of-record for review.

This architecture builds on the existing Physical AI trial infrastructure: the five-server MCP
architecture with 23 tools \cite{mcp-repo}, the federated learning framework with 235 modules
\cite{fl-repo}, the PSL and USL quantitative scoring standards \cite{national-platform-paper},
the 11-document site establishment package \cite{site-docs}, and the ten-stage patient journey
orchestration demonstrated across 168 patients in a 24-hour simulation \cite{main-repo}. The
national platform paper established the regulatory, standards, and infrastructure foundations;
the present paper completes the system by specifying the autonomous entity that sits atop this
infrastructure and executes the sponsor function.

The remainder of this paper is organized as follows. \S\ref{sec:governance} presents the
autonomous governance and portfolio management layer. \S\ref{sec:trial-design} describes the
trial design and statistical framework. \S\ref{sec:clinical-ops} covers clinical operations
and site management. \S\ref{sec:safety} details the safety and pharmacovigilance system.
\S\ref{sec:regulatory} addresses regulatory affairs and submissions. \S\ref{sec:quality}
presents quality assurance and compliance. \S\ref{sec:supply} describes supply chain and IMP
logistics. \S\ref{sec:data} covers data management and biostatistics.
\S\ref{sec:robotics} details the robotic execution gateway. \S\ref{sec:site-interface}
presents the site interface layer. \S\ref{sec:trust} describes the trust, provenance, and
audit infrastructure. \S\ref{sec:vendor} covers vendor management and CRO coordination.
\S\ref{sec:writing} addresses medical writing and disclosure. \S\ref{sec:financial} presents
the financial analysis. \S\ref{sec:implementation} outlines the national implementation
strategy. \S\ref{sec:discussion} provides a discussion of challenges and limitations.
\S\ref{sec:conclusion} concludes with a summary of contributions and future directions.
